\scanpage{215}
\mspage{1 ter}
% archive sheet number: 214

c'est prouvé pour les entiers $<n$. \{ Utilisons 3 \unsure{d},

\diagram{\begin{tikzcd}[ampersand replacement=\&,row sep=small,column sep=small]
X\arrow[d]\&X\times_Y X=X'\arrow[l]\arrow[d]\\
Y\&X=Y'\arrow[l,dash]
\end{tikzcd}}

et \unsure{regarder|regardons} le \deleted{\unsure{quotient|épimorphisme}}\inserted{morphisme} \unsure{diagonale|diagonal} $\Delta_{X/Y}$, \deleted{\unsure{par 2}} (existe par \textcircled{1}), par \textcircled{2} c'est un isom avec un sommand direct \inserted{universel}, donc $X'\simeq Y'\amalg X'_1$ ; par 3 a) $X'\longrightarrow Y'$

est de degré \unsure{$n$|$n-1$}, \deleted{\unsure{de}} par 3 \uline{c} $Y'\longrightarrow Y'$ est de degré 1, donc par 3 \uline{b} et \uline{f} on voit que $X'_1\longrightarrow Y'$ est de degré $n-1$. Lui appliquant l'hyp. de récurrence, on trouve $Y''_{\unsure{i|j|s}}\longrightarrow Y'$ \deleted{famille} de descente eff. universelle, qui trivialise $X'_1$ ; comme $X'\simeq Y'\amalg X'_1$ \uline{universellt},

on a $X''=X\times_Y Y''\simeq Y''\amalg\underbrace{(X'_1\times_{Y'}Y'')}_{Y''\times[1,n-1]}=Y''\times[1,n]$. Or $Y''\longrightarrow Y$ est de descente effective universelle, comme composé de deux tels, OK.

\uline{Explicitation} Pour vérifier qu'un \deleted{$f$} morphisme est de descente effective [cf condition 3 a)], il suffit de savoir que c'est un épimorphisme effectif universel, et de savoir que les relations d'équivalence sont effectives universelles \texttt{------}

\diagram{\begin{tikzcd}[ampersand replacement=\&,row sep=small,column sep=small]
{\tikz[baseline=-.3ex]{\draw[densely dotted,line width=.7pt] (0,.2)--(.17,.2)--(0,0)--(.17,0);}}\arrow[d]\&Z_X\arrow[l,dashed]\arrow[d]\&D\arrow[l,shift left=.5ex]\arrow[l,shift right=.5ex]\arrow[d]\\
Y\&X\arrow[l]\&X\times_Y X\arrow[l,shift left=.5ex]\arrow[l,shift right=.5ex]
\end{tikzcd}}

\noindent\rule{.25\linewidth}{.3pt}

\uline{Corollaire} Sous les conditions de A (iv), la notion de ``morphisme de degré $n$'' est uniquement déterminée, et \deleted{\unsure{corresp|correspond}} équivaut à : $\exists\,Y'\longrightarrow Y$ épimorphisme (\unsure{néces.} \deleted{\unsure{morphe}} effectif universel, i.e. de descente effec\ldots{}) avec $X\times_Y Y'\simeq[1,n]_{Y'}$ ($Y'$-isom).

